\chapter{Discussion}
\label{chap:discussion}

% See context/outline.md Ch 5 for paragraph-level plan.

This chapter interprets what \Cref{chap:findings} reported. It organises the primary findings under the three locked anchor concepts (Efficiency, Control, Adaptability), uses the layered theory built in \Cref{chap:theory} to read those findings, and then names the limitations that bound what the thesis can claim. \Cref{sec:primary-findings} carries the core argument under the anchors. \Cref{sec:adoption} draws out adoption implications. \Cref{sec:sustainability} works through sustainability and ethical considerations. \Cref{sec:profession-ethics} reflects on profession ethics, societal benefit, and the team. \Cref{sec:limitations} catalogues the twelve named limitations (L1 through L12) hierarchically. \Cref{sec:deviations} names deviations from the original plan. \Cref{sec:methodology-reflection} closes with a self-critical reflection on the central methodological weak spot.
% TODO-AUTO-ROUTER (lagt til 2026-05-10): Diskusjonskapitlet bør plassere
% auto-routeren (§3.4.5 i ch3) i akademisk sammenheng. Designet er strukturelt
% en portfolio-basert solver-velger (Gomes & Selman, *Algorithm Portfolios*,
% 2001; SATzilla som empirisk porteføljevelger for SAT-problemer) kombinert med
% en svært enkel utforsker-eksploiterer-tilnærming (slektskap til
% multi-armed bandit, jf. Sutton & Barto, *Reinforcement Learning: An
% Introduction*, kap. 2). Vår eksplor-strategi er forenklet (utforsker bare
% i cold start, slutter når N≥3), men er åpen for å erstattes med ε-greedy
% eller UCB i en senere iterasjon. Diskusjonen bør også kontrastere
% engine-valgs-tilpasning (system-styrt) med per-company weight-konfigurasjon
% (bruker-styrt) som to ulike Adaptability-leveranser.

This chapter interprets what \Cref{chap:findings} reported. It organises the primary findings under the three locked anchor concepts (Efficiency, Control, Adaptability), uses the layered theory built in \Cref{chap:theory} to read those findings, and then names the limitations that bound what the thesis can claim. \Cref{sec:primary-findings} carries the core argument under the anchors. \Cref{sec:adoption} draws out adoption implications. \Cref{sec:sustainability} works through sustainability and ethical considerations. \Cref{sec:limitations} catalogues the twelve named limitations (L1 through L12) hierarchically. \Cref{sec:deviations} names deviations from the original plan. \Cref{sec:methodology-reflection} closes with a self-critical reflection on the central methodological weak spot.


\section{Primary Findings under the Anchors}
\label{sec:primary-findings}

The discussion organises primary findings under the three locked anchor concepts: Efficiency, Control, and Adaptability, one sub-section per anchor, with the names used verbatim throughout.

\subsection{Efficiency}
\label{sec:efficiency}

The primary surprising result of \Cref{sec:interview-findings} is the resource-utilization visibility gap. Across the seven interviewed companies, overtime is handled reactively, idle time between assignments is invisible in current tools, and load balancing across drivers depends on coordinator memory rather than on any structured overview. The gap is not, on the face of it, a technological problem. The data needed to see utilization (driver hours, assignment durations, vehicle activity) exists in every company. What is missing is a structured view that makes utilization patterns legible to the people who would act on them.

The asymmetry that makes this gap a finding worth surfacing is the operator-vs-owner asymmetry. Owners and Admmit articulate utilization optimisation as something the company needs. Coordinators do not articulate it on their own. They ask for speed, override authority, and a system that does not get in their way. This is the configuration \textcite{bainbridge1983ironies} anticipates: demand for automation is articulated by those who design and own the system, while the residual operator role is left to whoever must run it day to day \parencite[p.~775]{bainbridge1983ironies}. Bainbridge's argument was theoretical when introduced in §2.2; the seven Norwegian transport interviews extend it empirically. The asymmetry is not an artefact of the interview sample: it appears across companies of different size, different current system, and different geographic location. The thesis does not claim the asymmetry is universal in Norwegian transport, but reports it as the consistent configuration in the consultation pool.

Three concrete utilization dimensions sit under \textbf{Efficiency} as the anchor names: overtime, idle time between consecutive assignments, and uneven load between drivers. RP makes each of them visible before any algorithm reorders them. The planning timeline shows assignments laid out across drivers and vehicles in a Gantt view; gaps between assignments are visible as gaps; long days with multiple consecutive assignments are visible as long bars; uneven load between drivers is visible as visual asymmetry across rows. The deviation viewer surfaces overtime violations as structured conflict data tied to specific assignments. The dashboard at \texttt{/oversikt} aggregates utilisation metrics across drivers and vehicles. Visibility itself is the precondition for optimization: a coordinator working from a spreadsheet or whiteboard cannot improve what they cannot see, and an algorithm cannot optimise what is not formalised.

The multi-engine benchmark in §4.5 is the validation instrument designed for Efficiency. It tests how solver approaches compare under realistic constraint combinations on identical synthetic problem instances; it does not test whether the visibility gap is real (the interviews surfaced this) or whether HITL is necessary (Admmit fixed this from project start). This how-not-whether framing is the methodologically independent test the thesis can defend. The headline result, traceable to commit \texttt{43c6124} on 2026-04-23, is the Timefold engine moving from being the worst-performing engine in its initial \texttt{HardSoftScore} configuration to being competitive with the heuristic baseline post-rework, an approximately seventy-two per cent reduction in canonical penalty across the three datasets. OR-Tools CP-SAT is the strongest individual engine on the instances it can finish (small and medium) but cannot complete the five-hundred-assignment large instance within the sixty-second time budget. That is the operational scaling boundary the thesis reports under Efficiency.

These findings indicate that RP delivers \textbf{Efficiency} along two distinct paths: visibility, which the seven coordinator interviews surfaced as missing in current practice and which the artefact addresses through the planning timeline, the deviation viewer, and the dashboard, and solver capability, which the multi-engine benchmark validates in Wieringa's how-not-whether sense rather than under deployment evidence.

\subsection{Control}
\label{sec:trust-control}

The argument that the coordinator should keep authority over every algorithm-generated assignment rests on Bainbridge's owner-vs-operator framing introduced in §2.2 and §5.1.1. Operator authority over override is not a usability feature added on top of an automated planner; it is the resolution to the asymmetry that motivated automation in the first place. \textcite{parasuraman2000automation} situate this configuration on their ten-level scale of automation for decision and action selection, where level 5 is "executes that suggestion if the human approves" and level 6 is "executes automatically unless the human vetoes within a time window". RP sits at level 5: the solver produces one plan, the plan is shown on the timeline, and no driver is dispatched until the coordinator approves it. Stopping deliberately at level 5 rather than going to level 6 makes override the constitutive interaction rather than the exception, which is what Control requires.

Whether a coordinator will trust the algorithm enough to use it is a separate question from whether they should keep authority. \textcite{hoff2015trust} review 127 studies and group their findings into a three-layer model of dispositional, situational, and learned trust \parencite[p.~413]{hoff2015trust}. Dispositional trust reflects how much an individual tends to trust automation in general. Situational trust depends on context and workload. Learned trust is built through experience with the specific system, first as initial trust based on reputation and then as dynamic trust that updates with each interaction \parencite[p.~420]{hoff2015trust}. \textcite{lee2004trust} sit beneath this three-layer model with the older claim that the design goal is appropriate trust, not greater trust: a coordinator who accepts every plan without review and a coordinator who reverts to fully manual planning despite a competent solver are both failures of trust calibration \parencite[p.~74]{lee2004trust}. Hoff and Bashir refines, rather than replaces, Lee and See: the three-layer model gives the developer a more granular set of antecedents to design against. Hoff and Bashir's finding that early errors damage trust more than later errors do \parencite[p.~426]{hoff2015trust} has a concrete consequence for RP: a coordinator's first week with the system weighs more in trust formation than the months that follow, and the deviation viewer's job in that first week is to make the algorithm's reasoning legible enough that an early misalignment can be inspected and corrected rather than read as a system fault.

Calibration depends on the coordinator being able to interpret what the algorithm has done, and on the system explaining itself in the coordinator's vocabulary rather than in solver internals. \textcite{miller2019explanation} argues that explanations are contrastive ("Why P rather than Q?"), selective (one or two relevant causes rather than a complete causal chain), and social (tailored to what the recipient already knows) \parencite[p.~3]{miller2019explanation}. RP shows contrastive reasons through the deviation viewer (each conflict flags type, affected assignment, and affected resource), through the post-solve summary panel (solution score and counts of hard and soft violations), and through the time-quality control surface introduced in §3.5.7 that labels the toggle as "fast suggestion" versus "best-effort plan" rather than as bare time labels.

The fourth element under Control is tacit knowledge as the coordinator's irreducible role. The interviews in §4.1 named tacit-knowledge dependency as the second most frequent pain point: every interviewed coordinator holds operational knowledge that is not formally captured by their current system, including who has the right competence on which routes, who is approaching overtime limits, and which vehicle suits which job. RP partially externalises this through the structured data the schema captures (competencies, certifications with expiry dates, work schedules, vehicle types and capacity) and through the deviation taxonomy that turns memorised constraints into structured conflict data on the timeline. What it cannot externalise is the coordinator's situated judgement: which route is harder in winter, which driver handles a difficult customer well, when to deviate from the algorithm because the day is unusual. The coordinator's review-and-override authority over each algorithm-generated assignment is how that irreducible judgement enters the system. Vague control language ("human oversight", "operator supervision") would not do the work that concrete review-and-override authority does.

These results highlight that meaningful \textbf{Control} requires four layers operating together: authority over plan content (override on every assignment), authority over plan-time process (the time-quality control), legible explanation in the operator's vocabulary (deviation viewer, post-solve summary panel, time-quality tooltip), and the structural decision to stop at Parasuraman level 5 rather than at level 6, so that approval becomes the constitutive interaction rather than the exception.

\subsection{Adaptability}
\label{sec:adaptability}

The seven interviewed companies operate under materially different rules. Bergen Bulk Transport has eight to twenty drivers and plans entirely manually. Norlog Tana runs around five routes per day across food-grade temperature-controlled transport with scanning competence requirements. Ottem operates forty-five vehicles across three departments with about a minute per assignment decision and unpredictable order cancellations. Nordic Crane runs crane and transport divisions with a planning horizon stretching from five hours to two months. Different fleet sizes, different operational rules, different willingness to pay. \textcite{venkatesh2003utaut} formalise this kind of variation as differences in adoption thresholds across organisational contexts; their UTAUT (Unified Theory of Acceptance and Use of Technology) model identifies four constructs (performance expectancy, effort expectancy, social influence, facilitating conditions) that shape technology adoption and use. The thesis does not test UTAUT empirically against the seven companies, but the model frames why a single artefact must accommodate genuinely different operational configurations to function meaningfully across them.

The frequency-ranked assignment-criteria list in §4.1 is the same across the seven companies (availability and position, experience, driving and rest-time status, competence and licence, route familiarity, equipment match, overtime status), but the relative weight each company places on each criterion differs. Some companies prioritise workload balance highly. Others prioritise vehicle--driver continuity or route familiarity. The variation is material rather than cosmetic: a plan that is good for one company can be operationally wrong for another, even when both face the same set of orders. The fit/gap analysis in §4.3 reinforces the point at the system level. The fourteen gap items are universal across the seven companies (planning view, automatic plan generation, conflict detection, capacity overview), but the five fit items (order management, billing, driver notification, basic data storage, status tracking) are met differently by Timpex, by Opter, by custom systems, or not at all, depending on the company.

Configurable soft-constraint weights are the technical mechanism that delivers Adaptability inside RP. Each of the five soft constraints (workload balance, vehicle preference, transitions, priority, employee preferences) has a per-company weight that flows into the CP-SAT objective function, so that the same engine produces different plans for different companies without code changes. The mechanism is the kind of variability point that \textcite{mietzner2009variability} describe in their multi-tenant SaaS variability framework: a single deployed instance carries a controlled set of configuration points that distinguish customers without forking the codebase. RP applies this to the soft-constraint layer of the planning problem: the weights are the per-tenant variability points; the engines, the data model, and the HITL interface are shared across tenants.

What is honest to acknowledge is that configurability for soft weights is a partial realisation of Adaptability. Hard constraints (HC-01 through HC-06) are hard-coded, the status workflow definition (planned, approved, in-progress, completed, cancelled) is hard-coded, and per-company assignment-rule taxonomies are hard-coded. A company whose hard constraints differ materially from the six implemented (for example, a company with regulatory rules outside Arbeidsmiljøloven, or a domain where double-booking is sometimes legitimate) cannot be served by configuration alone. Broader configurability of those layers is named as a gap in §5.5 and forwarded as future work in §6.3. Adaptability is also distinct from skalerbarhet, which concerns volume rather than meaningful operation across material rule differences. The thesis does not claim that RP scales to thousands of drivers; it claims that RP can serve materially different operational rulebooks within its current scope through configuration rather than code.

\textbf{Adaptability} is therefore demonstrated as a partial mechanism: sufficient for the soft-constraint layer through configurable weights and sufficient for the planning-horizon dimension, but insufficient for the hard-constraint and status-workflow layers that remain hard-coded. The gap between full and partial realisation is what the configurability work in \Cref{sec:future-work} carries forward.


\section{Adoption and Deployment Implications}
\label{sec:adoption}

What the artefact must demonstrate before companies adopt it is a cost-benefit threshold rather than a feature checklist. \textcite{venkatesh2003utaut} model adoption as the result of performance expectancy (will the system improve my work?), effort expectancy (will it be easy to use?), social influence (do people I respect use it?), and facilitating conditions (do the resources to use it exist?). The interviews in §4.1 mapped onto this directly: the sceptical interviewees (Harlem Solutions, Ottem, Bergen Bulk Transport) ranked cost-benefit and ease of use highest; the positive interviewees (the three Norlog branches) ranked performance and integration with existing systems highest. Bergen Bulk Transport's threshold for considering a TMS at all is fleet size: "I'll think about a TMS once we hit fifteen to twenty trucks." That is a performance-expectancy threshold expressed in fleet terms.

The TMS-as-category framing developed in §2.3 returns here as a deployment question. Existing systems (Timpex, Opter, internal tools) cover order management, invoicing, and driver notification. RP covers the planning step. They are complementary, not competing: a transport company that adopts RP keeps its existing TMS for invoicing. The most frequently named missing feature across interviews is integration with the existing billing system. Norlog Lakselv, Norlog Mo i Rana, and Harlem Solutions all flagged invoicing integration as an adoption-shaping requirement. The thesis acknowledges this gap honestly: billing integration is documented in §4.2 as FK-39, marked as Won't for this version, and recommended as future work in §6.3 \parencite{griffis2007tms, heinbach2022datadriven}.

What it would take to move from validation (this thesis) to evaluation (a deployed system used independently of the researchers) is the question §6.3 takes up. The minimum viable deployment is a production pilot with a real Norwegian transport company that bridges the gap between the synthetic-benchmark validation in §4.5 and Wieringa's evaluation criterion (\Cref{sec:dsr}). The pilot would address L5 (synthetic benchmarks, not production data) and L6 (no real-world deployment) at the same time, and would let the artefact's Efficiency claims be tested against operational utilization data rather than against synthetic instances. A second adoption barrier the interviews surfaced is performance: Timpex is described as extremely slow under concurrent use. RP's serverless deployment on a modern Next.js stack makes performance a differentiator rather than a parity feature, and the IFK-04 page-load target (under three seconds) and IFK-05 concurrent-user target (five to fifty coordinators per company) are deliberately chosen to clear the threshold the interviewees implicitly set.


\section{Sustainability and Ethical Considerations}
\label{sec:sustainability}

The Sustainability Awareness Framework (SusAF) developed by \textcite{becker2015karlskrona} and operationalised by \textcite{duboc2020requirements} structures sustainability analysis across five dimensions: technical, economic, social, individual, and environmental. The framework asks designers to consider both immediate effects and longer-term, indirect effects across each dimension. For an algorithm-assisted planning platform that intermediates the relationship between owners, coordinators, drivers, and customers, the framework is appropriate because the artefact's effects are not all in the same dimension and not all on the same time horizon.

The technical dimension covers the artefact's longevity: the multi-tenant architecture and the pluggable solver registry are the elements that support technical sustainability. New solvers can be added without touching the application; new companies can be added without touching the codebase. The economic dimension covers the artefact's value to the operating company: visible utilisation metrics support both internal cost control (overtime, vehicle utilisation) and external pricing decisions, and the per-company configurability supports a multi-tenant deployment model that distributes development cost across customers. The social dimension covers the artefact's effects on the people working with and around it: the HITL configuration deliberately keeps the coordinator's judgement in the loop rather than replacing it, but the system's optimisation of workload balance can also be read as tightening the workload baseline a driver is expected to meet. The individual dimension covers the coordinator's professional identity: a planner whose tacit knowledge has been the source of their value can experience formalisation as a deskilling threat, even when the formalisation was meant to remove key-person risk. The environmental dimension covers the artefact's effects on physical resource use: better utilisation reduces idle time and unnecessary trips, which translates to lower fuel use and emissions, but the magnitude of that effect on the seven interviewed companies is not measured by the thesis and cannot be claimed without further empirical work.

\Cref{tab:susaf} maps the analysis above onto the SusAF five-dimension grid. Cells marked with `---' are dimensions where this thesis does not have empirical or analytical material to make a claim, rather than dimensions where no effect exists; closing those gaps is part of what production deployment in \Cref{sec:considering-deployment} would enable.

\begin{table}[h]
\centering
\caption{Sustainability Awareness Framework (SusAF) analysis of RP across five dimensions and three orders of effect (immediate, enabling, structural). Empty cells mark places this thesis does not have empirical or analytical grounding to fill, not places where no effect exists.}
\label{tab:susaf}
\small
\begin{tabular}{p{0.13\textwidth}p{0.27\textwidth}p{0.27\textwidth}p{0.23\textwidth}}
\toprule
\textbf{Dimension} & \textbf{Immediate} & \textbf{Enabling} & \textbf{Structural} \\
\midrule
Technical      & Multi-tenant architecture and pluggable solver registry support longevity         & New solvers and new companies can be added without touching the application code & --- \\
Economic       & Visible utilisation metrics support internal cost control and pricing decisions   & Per-company configurability supports a multi-tenant SaaS deployment model        & Distributes development cost across customers \\
Social         & HITL keeps the coordinator's judgement in the loop rather than replacing it       & Workload-balance optimisation can tighten the baseline drivers are expected to meet & --- \\
Individual     & Tacit-knowledge formalisation can be experienced as a deskilling threat           & ---                                                                              & --- \\
Environmental  & Better utilisation reduces idle time and unnecessary trips $\rightarrow$ lower fuel use and emissions & --- (magnitude not measured)                                              & --- \\
\bottomrule
\end{tabular}
\end{table}

Four ethical dilemmas tie sustainability to specific design questions. Fairness, in the sense \textcite{mittelstadt2016algorithms} survey for algorithmic decision-making, becomes a question about which drivers the algorithm prefers and why: a workload-balance term that smooths total hours can systematically advantage or disadvantage individual drivers depending on the priority weighting and the assignment distribution. Accountability, as \textcite{martin2019accountability} argue, is the question of who is answerable when an algorithm-generated assignment goes wrong, and the answer in RP ties directly to \textbf{Control}: the coordinator who reviewed and approved (or overrode) the assignment is the answerable party, not the algorithm or the company that built it. Privacy is the question of how employee data (competencies, work schedules, time-off, sick-leave, performance) is handled inside the system; RP keeps this data scoped to the employing company through multi-tenant isolation, but a richer feature set (driver-facing app, GPS tracking, real-time location) would push privacy into a different regime. Working conditions tie to \textbf{Efficiency}: \textcite{lee2018understanding} document how algorithmic management can compress task slack and tighten worker discretion in ways that erode the qualitative experience of work, and the question for RP is whether visibility into utilisation drives systemic optimisation that benefits the company at the cost of the drivers being optimised over.

Mapping these effects onto the United Nations 2030 Agenda \parencite{un2015agenda2030} is illustrative rather than exhaustive. SDG 8 (Decent work and economic growth) is the dimension where Efficiency improvements (lower overtime, more even workload) and the working-conditions dilemma cut in opposite directions. SDG 9 (Industry, innovation and infrastructure) is where the multi-tenant SaaS deployment model and the open-source solver components fit. SDG 13 (Climate action) is where lower idle time and fewer unnecessary trips are relevant, with the caveat that the magnitude is not measured.

Two more notes. The operator-vs-owner asymmetry that motivates the artefact is itself an ethics question: whose efficiency, whose autonomy? The thesis comes down on the side that lets coordinators keep authority over individual assignments while making utilisation visible to owners, but a different design that pushed authority upward (an automatically-deployed plan that coordinators can only veto in time-critical windows) would be a different ethical commitment. And the sustainability analysis here is analytical, not empirical: it identifies effects the artefact could plausibly have, but does not measure them on the seven interviewed companies. That bound is named as part of L5 and L6 in §5.5.


\section{Ethical Reflections for the Profession - Societal Benefits and Impact}
\label{sec:profession-ethics}

The sustainability analysis in \Cref{sec:sustainability} looks outward, at the artefact's effects on people and systems beyond the project. This section turns the same question inward: the responsibilities the work carried as a piece of professional engineering, what it plausibly delivers to the people it was built for, and how the team functioned through the year.

Three commitments shaped the execution. The first is anonymisation. The seven traffic coordinators and transport managers interviewed are referenced by company-anonymised codes throughout, and the consent flow established in \Cref{sec:data-collection} is what gives the empirical claims their footing. The second is honesty about validation status. The multi-engine benchmark in \Cref{sec:evaluation-framework} runs on synthetic datasets, the artefact has not been deployed in production, and no operating coordinator has used it independently of the team. The thesis names these bounds as L5, L6, and L8 in \Cref{sec:limitations} rather than soft-pedalling them, because a thesis whose limits are camouflaged is a thesis that cannot be built on. The third is due credit. \Cref{sec:recommendations} addresses what the project recommends to Admmit specifically, and the preface names the people whose support made the work possible.

The societal benefit the work plausibly delivers is concentrated on the working day of the traffic coordinator. The visibility gap reported in \Cref{sec:primary-findings} is felt every shift: overtime emerges reactively, idle time is invisible, and load balancing relies on memory. An artefact that surfaces utilisation as the plan is built, leaves the coordinator's authority over each assignment intact, and explains its reasoning in the operator's vocabulary lowers the cognitive load of work currently held together by spreadsheets and recall. Three other groups stand to benefit indirectly: drivers, through more even workload distribution and fewer reactive overtime calls; owners, through cost control grounded in measured utilisation rather than inferred from invoices; and Admmit's customers more broadly, through software that reflects how Norwegian transport actually plans rather than how a generic optimisation engine assumes it should be planned.

The team worked as a pair throughout the year, with division of labour decided iteration by iteration rather than fixed up front. The pair structure was a strength for review, since every artefact and every chapter had a second pair of eyes before submission, and a constraint for breadth, since seven interviews and three solver engines stretched a two-person team thinly. Individual reflection notes are submitted alongside this thesis.


\section{Limitations}
\label{sec:limitations}

The L1 through L12 names below are locked verbatim from the limitations register; synonyms (e.g., "smaller sample" for L1) drift the spine and must be flagged by reviewers. \Cref{tab:limitations} summarises the register; each entry is then expanded across the three sub-sections that follow.

\begin{table}[h]
\centering
\caption{Hierarchical limitations register. Scope categories group the limitations into empirical foundation, validation and artefact, and conceptual or methodological bounds; the SQ column records which sub-question each limitation principally bounds.}
\label{tab:limitations}
\small
\begin{tabular}{p{0.05\textwidth}p{0.18\textwidth}p{0.62\textwidth}p{0.05\textwidth}}
\toprule
\textbf{ID} & \textbf{Scope} & \textbf{Statement} & \textbf{SQ} \\
\midrule
L1  & Empirical          & Sample size of seven brief interviews; thematic coverage rather than statistical generalisation & SQ1 \\
L2  & Empirical          & Self-selection bias: recruitment from Admmit's customer roster only                              & SQ1 \\
L3  & Empirical          & Researcher--developer overlap; informal study procedures under consultation framing              & SQ1 \\
L4  & Empirical          & Interview-guide coverage gap: rule differences across companies not exhaustively probed           & SQ1 \\
\midrule
L5  & Validation         & Synthetic benchmarks rather than production data                                                  & SQ3 \\
L6  & Validation         & No real-world deployment of RP                                                                    & SQ3 \\
L7  & Validation         & No empirical comparison against existing TMS in operational use                                   & SQ3 \\
L8  & Validation         & No user testing of the override flow with traffic coordinators                                    & SQ2 \\
L9  & Validation         & Algorithm evaluation against the project's own benchmarks only                                    & SQ3 \\
\midrule
L10 & Conceptual         & HITL as Admmit mandate, not as a validated automation level                                       & SQ2 \\
L11 & Conceptual         & Single-domain (Norwegian transport)                                                               & SQ3 \\
L12 & Conceptual         & Boundary cases described qualitatively, not quantified                                            & SQ3 \\
\bottomrule
\end{tabular}
\end{table}

\subsection{Empirical Foundation}
\label{sec:limits-empirical}

\paragraph{L1 — Sample size (7 interviews).} Seven semi-structured calls of approximately fifteen minutes each is a small empirical foundation. The sample was chosen for thematic coverage rather than for statistical generalisation, and the analysis in §3.4 explicitly treats the themes as bounded to the seven interviewed companies rather than as a basis for population-level claims. What L1 prevents the thesis from claiming is anything quantitative about Norwegian transport at scale: it cannot say that the visibility gap holds in eighty per cent of Norwegian transport companies, only that it surfaced consistently across the seven calls in the consultation pool. The brief duration compounds this: the calls were long enough to surface the visibility gap and the operator-vs-owner asymmetry but not long enough to probe edge cases or to test whether a coordinator's articulated need shifts under sustained reflection.

\paragraph{L2 — Self-selection bias in interview sample (Admmit customers only).} The seven coordinators were recruited through Admmit's customer roster rather than through a public call. Admmit's customers are companies that have already chosen to engage with a software vendor in the planning space, which is a self-selected slice of Norwegian transport. What L2 prevents the thesis from claiming is anything about transport companies that have explicitly chosen not to engage with planning software, including a substantial share of small operators that, per Bergen Bulk Transport's own articulation, will not consider a TMS until the fleet exceeds fifteen to twenty trucks. The operator-vs-owner asymmetry in particular may be sharper inside Admmit's customer base than across the broader population, because the asymmetry is part of why owners brought Admmit in.

\paragraph{L3 — Researcher–developer overlap and informal study procedures.} The two authors who conducted and analysed the interviews are the same two people who designed and built the artefact, and the interviews were conducted as industry consultation under Admmit's bachelor task brief rather than as formal research, with the procedural consequences set out in §3.3 (no formal consent form, no formal anonymisation protocol). \textcite{malterud2001lancet} treats preconceptions as separate from bias unless the researcher fails to disclose them \parencite[p.~484]{malterud2001lancet}; the overlap and the informal procedures are declared in §3.7 and named here as a limitation that bounds reflexivity. What L3 prevents the thesis from claiming is full independence between the empirical work and the design commitments: the prior expectation that a visibility gap might exist, and the prior commitment to HITL that the bachelor task fixed in advance, both predate the interviews and were carried into them by the same people who would later build the artefact. The overlap does not invalidate the empirical findings, but it does mean that the asymmetry between Admmit's articulation and the coordinators' articulation should be read as a triangulation point, not as an unmediated discovery.

\paragraph{L4 — Interview-guide coverage gap.} The five guide topics (current tools and workflow, assignment criteria, sick-leave handling, attitudes to automation, adoption criteria) covered the questions the project most needed answered, but they did not exhaustively probe what would belong in a richer Adaptability evaluation. Specifically, the guide did not ask coordinators to enumerate the hard constraints they consider non-negotiable beyond licence and competence, did not probe the status workflow definitions across companies, and did not ask about per-company assignment-rule taxonomies. What L4 prevents the thesis from claiming is that the configurability gap named in §5.1.3 (hard constraints, status workflow, assignment rules remain hard-coded) is small. A richer interview guide might have surfaced rule differences large enough to reposition the gap as central rather than peripheral.

\subsection{Validation and Artefact}
\label{sec:limits-validation}

\paragraph{L5 — Synthetic benchmarks, not production data.} The benchmark in §4.5 runs against three synthetic datasets generated to span the fleet ranges named in §3.3. Synthetic data has known properties (controlled overlap density, controlled constraint intensity) that production data does not have. What L5 prevents the thesis from claiming is anything about how RP would behave on real Norwegian-transport operational data. The synthetic generator may understate the constraint pressure of a real day and overstate it on others; the patterns of workload balance and route familiarity that real coordinators would recognise are not present. The benchmark is validation evidence in Wieringa's sense (\Cref{sec:dsr}), not evaluation evidence: it predicts how the engines compare under controlled inputs, not how they perform when stakeholders use the artefact in production.

\paragraph{L6 — No real-world deployment.} No transport company has operated RP independently of the researchers. The artefact is implemented and developer-verified; thirty-six of the thirty-seven in-scope functional requirements pass the implementation-and-developer-verification check in the requirements traceability matrix in §3.6. What L6 prevents the thesis from claiming is that the system delivers measurable Efficiency gains under operational time pressure, that the override flow holds up under coordinator workload, or that adoption barriers identified in §5.2 are correctly weighted. \textcite{hevner2007threecycle} treats the lab-before-field ordering as expected DSR practice \parencite[p.~91]{hevner2007threecycle}, so L6 is a bound on the project's stage rather than a methodological deficit; it is, however, the limitation most central to the thesis's external validity.

\paragraph{L7 — No empirical comparison against existing TMS.} The fit/gap analysis in §4.3 compares RP's coverage against what existing systems cover, but it does not benchmark RP against Timpex, Opter, or the internal and custom tools the seven interviewed companies use. What L7 prevents the thesis from claiming is that RP is better than existing alternatives in the same market segment. The differentiator the thesis can defend is structural (RP fills the planning gap that existing systems do not address), not comparative (RP outperforms alternatives on a shared task). Closing this gap would require operational deployment alongside an existing system, which L6 also bounds.

\paragraph{L8 — No user testing with coordinators.} The HITL interface (the planning timeline, the deviation viewer, the time/quality control, the drag-and-drop override) is implemented and developer-verified, but it has not been user-tested with traffic coordinators independently of the development team. What L8 prevents the thesis from claiming is that the override flow is comprehensible, that the time/quality labels work for coordinators who do not think in solver-runtime budgets, or that the deviation taxonomy surfaces conflicts at a useful threshold. The Control argument in §5.1.2 is design argument plus theory; user testing would convert it to design argument plus theory plus empirical evidence. L8 is the limitation that most directly bounds SQ2.

\paragraph{L9 — Algorithm evaluation against own benchmarks only.} The benchmark in §4.5 reports results against three synthetic datasets that the project itself defined. There is no independent benchmark suite that RP's engines run against, and no public dataset standard for the assignment problem RP solves. What L9 prevents the thesis from claiming is that the engine performance numbers transfer to other formulations of related problems. The Timefold seventy-two per cent canonical-penalty reduction is a real, measurable engineering result against the canonical scorer in this codebase, but cross-codebase comparison would require a shared dataset and a shared scorer that the project does not currently have.

\subsection{Conceptual and Methodological}
\label{sec:limits-conceptual}

\paragraph{L10 — HITL as mandate, not validated level.} HITL was Admmit's requirement from project start. The seven interviews validated the necessity of HITL through coordinator consensus that automated planning must be correctable rather than authoritative, but they did not validate the specific Parasuraman level (level 5: system suggests, human approves) at which RP sits. What L10 prevents the thesis from claiming is that level 5 is the right level for Norwegian transport rather than, say, level 4 (system suggests an alternative, human selects) or level 6 (system executes automatically unless vetoed within a time window). A coordinator who actually used RP under operational time pressure might prefer level 6 for routine assignments and level 5 for exceptional ones; the thesis cannot tell.

\paragraph{L11 — Single-domain (Norwegian transport).} The empirical foundation, the artefact, and the discussion are all anchored in Norwegian transport. What L11 prevents the thesis from claiming is that the visibility gap, the operator-vs-owner asymmetry, the multi-engine architecture, or the configurability mechanism transfer to adjacent domains (other Nordic transport markets, other personnel-scheduling settings such as nurse rostering or airline crew scheduling) without further empirical work. Cross-domain replication would address L11 directly and is named as future work in §6.3.

\paragraph{L12 — Boundary cases described, not quantified.} Several boundary cases in §3.5 (false-positive deviation alerts at week boundaries and mid-shift absences, pathological soft-constraint weight combinations, the lex contract masking Timefold's superiority on the legal and business-critical tiers on the large instance) are described as observations from the iterative development process. They are not quantified across a sample of cases, and the thesis does not report frequency distributions for any of them. What L12 prevents the thesis from claiming is a precise rate at which any boundary case occurs in normal operation. Operational deployment would let these be measured; until then, they remain qualitative observations carried as caveats on the engine results in §4.5.


\section{Deviations from Plan}
\label{sec:deviations}

Three deviations from the original plan are worth naming honestly. First, the seven interviews were originally scheduled to be conducted in a single day. They were instead spread across the first half of March 2026, with one call at a time and reflection between calls so that gaps in earlier conversations could be probed in later ones. The deviation produced a stronger empirical foundation than the single-day plan would have: the visibility gap and the operator-vs-owner asymmetry both surfaced through cross-call comparison rather than from any single conversation. Second, the multi-engine architecture was not in the original scope. The greedy baseline in §3.5.2 was the planned algorithm; OR-Tools CP-SAT was added in §3.5.3 after the greedy heuristic's locally-good-globally-suboptimal pattern surfaced; Timefold was added in §3.5.4 after the OR-Tools scaling boundary surfaced. The deviation produced the methodologically independent how-not-whether benchmark in §4.5 and shifted the thesis's central evaluation from "the system is feasible" to "the system's solver layer is comparable across families". Third, user testing with coordinators was on the original work plan and was not delivered. The decision to forgo user testing was driven by the project's stage (§3.6 frames the work as validation in Wieringa's sense rather than evaluation) and by the limited time available after the eight implementation iterations stabilised. The deviation is named as L8 in §5.5 and as future work in §6.3, and is the most consequential of the three for what the thesis can claim about Control.


\section{Methodology Reflection}
\label{sec:methodology-reflection}

This section reflects on the research methodology in three named pillars: reproducibility (whether the work can be repeated), validity (whether the study examines what it claims to examine), and setbacks (what went wrong during construction and how the artefact and the evaluation were shaped by it).

\subsection{Reproducibility}
\label{sec:reproducibility}

What supports reproducibility in this thesis is a deliberate set of choices made during the multi-engine benchmark's design. Engine versions are snapshotted per benchmark run because early benchmark inconsistency in §3.5.4 traced to silent solver-library updates between runs; the snapshotting removed that source of drift. The three engines (greedy, OR-Tools CP-SAT, Timefold) read the same shared constraint configuration loaded from a single file, so cross-engine comparisons measure search behaviour rather than constraint definitions. Each returned solution is re-scored using one shared scorer with lex-tier scoring (Norwegian Arbeidsmiljøloven legal constraints, then hard constraints, then soft constraints) so engine-specific scoring conventions do not contaminate comparison. The headline Timefold reduction from §4.5 is traceable to commit \texttt{43c6124} on 2026-04-23, a precision the engineering result depends on.

What limits reproducibility is named in §5.5. There is no public dataset standard for the assignment problem RP solves (L9), so cross-codebase comparison would require a shared dataset and a shared scorer the project does not currently provide. There is no automated test suite for non-functional behaviour (also part of L9), so the requirements traceability matrix's "developer-verified" claim is narrower than passing a test suite would be. Wall-clock figures are single-sample (§3.6.1), and engine-internal randomness is therefore part of the variation a single sample captures rather than something averaged out across multiple trials.

\subsection{Validity}
\label{sec:validity-reflection}

The most load-bearing methodological weak spot in this thesis is the synthetic-benchmark validation that carries Efficiency. The benchmark is the central evaluation instrument for the artefact's primary anchor. It is rigorous within its own terms, but what it cannot do is tell the reader how RP behaves on real Norwegian-transport operational data. The synthetic generator was tuned to span the fleet ranges named in §3.3, but synthetic overlap density, synthetic priority distributions, and synthetic time windows are not the same as real ones, and a coordinator using RP on a real day would face constraint pressure the synthetic datasets may understate or overstate. L5 names this bound, L6 names the absence of deployment that would close it, and L7 names the absence of comparison against existing systems that would let the differentiator the thesis claims (planning gap coverage) be tested empirically. If the project had a second year, the most valuable methodological move would be a production pilot with one of the seven interviewed companies, run against a real day's operational data with the coordinator's own override decisions captured. Without that, the thesis's Efficiency claim is structurally a claim about how solvers compare under controlled inputs, not a claim about utilisation gains in deployed Norwegian transport. The thesis is honest about that distinction throughout, but the distinction is the most significant single qualifier on what the work establishes.

\subsection{Setbacks}
\label{sec:setbacks}

Three setbacks during construction shaped the artefact and the evaluation in ways worth naming explicitly, alongside the deviations already catalogued in §5.6.

First, the Timefold engine ran for several weeks under an initial \texttt{HardSoftScore} configuration that proved insufficient for the assignment problem's tiered-priority structure. The constraint-provider rework on 2026-04-23 corrected this and is the source of the seventy-two per cent canonical-penalty reduction reported in §4.5, but several intermediate iteration outcomes were therefore measured against an unfair Timefold baseline before the rework landed.

Second, the lex-tier contract built into the auto portfolio ranks engines by hard-tier penalty first, which on the large instance masks the case where Timefold beats greedy on the legal, business-critical, and preference tiers but loses on hard. The contract was kept rather than tuned because the operational consequence of a single hard-tier violation outweighs gains on lower tiers, but the masking is a methodological artefact worth naming.

Third, the post-hoc deviation taxonomy produced false-positive alerts at week boundaries and mid-shift absences during early HITL iterations. The boundaries between the weekly work-schedule definition and time-off windows had to be reconciled in the schema before the deviation flags became reliable, and the rest-period regulation has multiple definitions (the forty-eight-hour weekly average, the sixty-hour weekly maximum, the daily thirteen-hour limit) that produce inconsistent flags depending on which definition the runtime check applies.

None of the three setbacks is a flaw the project failed to address; each is a step the artefact passed through and a methodological artefact the thesis carries as caveats on its results.


\section{Recommendations for Admmit and the Norwegian Transport Software Market}
\label{sec:recommendations}

The thesis's empirical foundation and construction experience yield four specific recommendations. The first three are addressed to Admmit as the industry partner that brought the bachelor task; the fourth is addressed to the Norwegian transport software market more broadly. None substitutes for the deployment work named in §5.5 as the central absence; each is a project-experience-based recommendation that does not depend on deployment evidence to act on.

\paragraph{Open the configuration surface beyond soft-constraint weights.} Adaptability is currently delivered through per-company soft-constraint weights only. The configurability gap in §5.1.3 names hard constraints (HC-01 through HC-06), the status workflow definition (planned, approved, in-progress, completed, cancelled), and per-company assignment-rule taxonomies as the layers that remain hard-coded. A subsequent product iteration that exposes these to per-tenant configuration, with an interview-guide pass that probes the rule differences L4 names, would convert Adaptability from partial mechanism into a full one and reduce the integration cost of onboarding a company whose hard constraints differ materially from the six implemented.

\paragraph{Invest in the existing-TMS integration path rather than in a competing billing module.} The fit/gap analysis in §4.3 established that Timpex and Opter cover order management, invoicing, and driver notification, and that RP covers the planning step neither addresses. The most frequently named missing feature across interviews was integration with the existing billing system. Building a competing billing module would put RP into direct competition with the systems it depends on for adjacent functions; building a documented integration path keeps RP in the planning-and-decision-support category the fit/gap analysis identified as its differentiator. The integration project is bridging rather than replacing.

\paragraph{Front-load the operator first-week experience.} \textcite{hoff2015trust} report that early errors in human-automation interaction damage trust more than later errors do \parencite[p.~426]{hoff2015trust}. For RP this means a coordinator's first week with the system weighs more in trust formation than the months that follow. Concrete implications for product rollout: structured onboarding rather than self-serve, supervisor proximity for deviation-flag interpretation, the deviation viewer's explanations tested for legibility before the first coordinator runs an operational plan, and reluctance to ship features that would expose the coordinator to early misalignments before the system has earned the calibrated trust §5.1.2 sets out as the design target.

\paragraph{Adopt or push for a shared synthetic-dataset format for Norwegian transport.} L9 in §5.5 names the absence of a public benchmark suite for the assignment problem RP solves; cross-codebase comparison currently requires a shared dataset and a shared scorer the project does not provide. A small standard test-data format agreed across Admmit, Timpex, Opter, and the internal-tool builders would let multi-vendor benchmarking become possible, would let the kinds of comparison the thesis cannot perform under L7 be performed in subsequent work, and would make the Norwegian transport software market comparable in a way the order-management category already is and the planning category currently is not.


\section{Considering Production Deployment}
\label{sec:considering-deployment}

The thesis validates RP in Wieringa's sense (predicting how a deployed system would behave) rather than evaluating it in deployment. \Cref{sec:dsr} sets out the validation-versus-evaluation distinction; this section asks what concrete steps would move RP from validation to evaluation. Three sub-sections name what each of the three core limitations (L5, L7, L8) would require to close, in operational rather than methodological terms.

\subsection{From Synthetic to Operational Data (L5)}
\label{sec:deployment-data}

What L5 would require to close is a real day's order pool from one of the seven interviewed companies, with anonymised driver, vehicle, and competency data, run through the same multi-engine benchmark protocol the thesis applies to its synthetic datasets, and with the coordinator's override decisions captured as the ground truth against which engine output is compared. The seven coordinator interviews supply candidate pilot companies. Norlog Mo i Rana, which described Timpex performance as the dominant pain point, is a candidate where RP's performance would visibly differentiate against current practice. Bergen Bulk Transport, the smallest interviewed company, is a candidate where the constraint pressure on the small-instance benchmark could be cross-checked against a real day. Closing L5 is therefore not a research question; it is an operational data-collection task whose ethical envelope has been set by the consultation framing the thesis already operates inside (§3.3).

\subsection{From Validation Against Self to Comparison Against Alternatives (L7)}
\label{sec:deployment-comparison}

What L7 would require to close is a parallel-running configuration in which the same coordinator generates the same day's plan with their current tool (Timpex, Opter, an internal tool, or manual) and with RP, and the two plans are compared on the shared lex-tier scorer §4.5 establishes. The comparison the thesis cannot make today is not a head-to-head feature comparison (the fit/gap analysis in §4.3 establishes that RP and existing systems sit in different categories); it is a utilization-outcome comparison on the dimensions Efficiency names. Whether RP's planning-gap closure produces measurable reductions in overtime, idle time, and uneven driver load relative to whatever each company does today is the comparison L7 leaves open and that a parallel-running configuration would close.

\subsection{From Override-Flow Specification to Coordinator User Testing (L8)}
\label{sec:deployment-usertest}

What L8 would require to close is a coordinator using RP on actual orders with the system recording each override event: which assignment was overridden, what the algorithm originally suggested, what the coordinator changed it to, and any reason the coordinator chose to record. Seven days of such logging across two or three coordinators would supply the data the Control argument in §5.1.2 currently asserts on theoretical grounds and would let the SQ2 answer move from design plus theory to design plus theory plus operational evidence. The override-event log is also the instrument through which the system can detect when a coordinator's first-week trust calibration is going wrong, which §5.8's first-week recommendation depends on observability for.

\subsection{The First Pilot}
\label{sec:first-pilot}

Closing all three at once is not the right starting point. The lowest-friction first move is L8 (override-flow user testing) at one of the smaller interviewed companies, which generates the operational override-event data the Control argument needs and surfaces the implementation gaps the synthetic benchmark cannot. L5 and L7 follow naturally once one company is running RP day-to-day: the order pool is then real, and the comparison against the company's previous tool becomes available to score. The pilot is named in \Cref{sec:future-work} as the highest-priority follow-on work, and the structural argument the thesis closes on (algorithm-assisted planning under stakeholder asymmetry is feasible when authority lives with the operator, explanation lives in the operator's vocabulary, and configuration lives between companies rather than inside the code) becomes empirically testable through it.
